\section{Legends and Later Images}

The historical Paul ends at the Ostian grave; the remembered Paul was just
beginning. The legends deserve a chapter of their own---not to be
believed, but because they carried his image for centuries and still
supply, unlabeled, much of what people think they know about him.

\subsection{The \work{Acts of Paul} and Thecla}

The great source is the apocryphal \work{Acts of Paul}, a novelistic cycle
that Tertullian says was ``composed shortly before his time in honour of
Paul by a presbyter of Asia, who was convicted of the imposture and
degraded from his office''---hence datable c.~160, an ``orthodox
Christian'' author by M.~R. James's judgment, whose 1924 translation is
used here. Its most famous episode is Thecla of Iconium: a betrothed
virgin who ``sat at the window hard by, and hearkened night and day unto
the word concerning chastity which was spoken by Paul,'' broke her
engagement, was condemned to the fire and then to the beasts, baptized
herself in the arena tank---``In the name of Jesus Christ do I baptize
myself on the last day''---was spared, and was sent out by Paul with the
commission ``Go, and teach the word of God.'' The text closes her life
where her later cult flourished: ``she departed unto Seleucia, and after
she had enlightened many with the word of God, she slept a good sleep.''
Thecla became one of the most beloved saints of the early church; her
story is second-century romance, not first-century record, and its
picture of Paul as preacher of radical celibacy is the author's ascetic
program wearing the apostle's cloak---but it is also the earliest surviving
evidence of how powerfully Paul's memory could authorize a woman's
ministry in some second-century circles.

\subsection{The face of Paul}

The same Thecla narrative opens with the only ancient physical description
of the apostle, given to Onesiphorus, who waits by the highway for a man
he has never seen: ``And he saw Paul coming, a man little of stature,
thin-haired upon the head, crooked in the legs, of good state of body,
with eyebrows joining, and nose somewhat hooked, full of grace: for
sometimes he appeared like a man, and sometimes he had the face of an
angel.'' This is literature, not portraiture---a century removed from its
subject, and possibly a rhetorical ideal (the description is not
unflattering by ancient physiognomic conventions). But it is the ancestor
of the entire iconographic type---the small, bald, bearded,
beetle-browed Paul of the catacomb frescoes, mosaics, and icons, with the
sword of his martyrdom as attribute---and it may embroider a genuine
memory; Paul's opponents at Corinth, after all, said ``his bodily
presence is weak'' (2~Cor 10:10). Portraits were in circulation early:
Eusebius, writing before 340, attests that ``the likenesses of his
apostles Paul and Peter, and of Christ himself, are preserved in
paintings,'' a practice he attributes to converted Gentile habits of
honoring benefactors (\work{Church History} 7.18). Iconography records
reception, not appearance. Status of the description: L, early legend
possibly preserving a physical reminiscence, unverifiable.

\subsection{The martyrdom legend: milk, and the preaching corpse}

The \work{Acts of Paul} also supplies the earliest martyrdom narrative,
and its flavor should be tasted. Nero, alarmed by the Christians in his
own household, ``commanded all the prisoners to be burned with fire, but
Paul to be beheaded after the law of the Romans.'' At the block, ``Paul
stood with his face to the east and lifted up his hands unto heaven and
prayed a long time\ldots\ and then stretched forth his neck without
speaking. And when the executioner (speculator) struck off his head, milk
spurted upon the cloak of the soldier.'' The dead apostle then appears to
Nero---``Caesar, behold, I, Paul, the soldier of God, am not dead, but
live in my God''---and to the believing officers at his grave. The milk
(purity, or the nurse-image Paul used of himself) is symbol, not
coroner's report; the narrative's value is double: it fixes beheading as
the received mode by c.~160, and it shows the legend already outrunning
the history. The fifth-century \work{Acts of Peter and Paul} of
Pseudo-Marcellus added the local etiology: the severed head bounded three
times, and at each touch a spring broke out---the \textit{Tre Fontane},
``three fountains,'' at the Aquae Salviae, where a monastery and church
stand to this day. Benedict~XVI's catechesis labels this element frankly
``the most legendary'' of the traditions (4 February 2009); the site
itself, as the traditional scene of the execution, is venerable but
unverifiable.

\subsection{Paul and Seneca}

By the fourth century cultured Christians could read a correspondence
between Paul and his exact contemporary Seneca---brother of the very
Gallio of Acts 18. Jerome, who wanted the Stoic on his shelf, admits the
letters got Seneca into his catalogue of Christian writers: ``I should
not place him in the category of saints were it not that those Epistles
of Paul to Seneca and Seneca to Paul, which are read by many, provoke
me'' (\work{De viris illustribus}~12). The fourteen Latin letters are a
late, artless forgery, doubted since the Renaissance and defended by no
one; their interest is what they witness: the hunger, already ancient, to
seat the tent-maker among the philosophers.

\subsection{Grading the strata}

The pattern across all these traditions is regular and worth stating as
method: the closer to the man, the sparer the record; the further away,
the richer. First-century sources give no portrait, no execution scene,
no springs, no philosophical correspondence; the second century supplies
face and martyrdom narrative; the fourth and fifth supply topography and
famous friends. Each layer is evidence---for its own age's devotion, for
the cult's geography, for iconography---and none of it may be passed
backward as biography. The principal items, graded:

\begin{fourstable}{Tradition}{Earliest located witness}{Later
reception}{Status}
\properrow{Martyrdom at Rome under Nero}{1 Clem.\ 5 (c.~95--96)}
{universal; no rival claim}{early, convergent; morally certain (E/S)}
\properrow{Death by beheading}{\work{Acts of Paul} (c.~160); Tertullian,
\work{Praescr.}\ 36 (c.~200)}{universal; the sword as attribute}{early
tradition, not contemporary; coheres with citizenship (E/L)}
\properrow{Same day as Peter}{Jerome, \work{De viris} 5 (392); Dionysius
c.~170 says only ``at the same time''}{joint feast, 29 June}{calendar
piety; not chronology (E/L)}
\properrow{Journey to Spain}{Rom 15:24 (intention); Muratorian fragment
(assumes it); 1 Clem.\ 5 (ambiguous ``limit of the west'')}{release
tradition via Eusebius, Jerome}{intention certain; execution unproven
(A/E)}
\properrow{Ostian burial; Lucina's estate}{Gaius c.~200 (the trophy);
Pseudo-Abdias (4th c.) for Lucina}{basilica; continuous shrine}{site
tradition early and unrivaled; detail legendary (E/L/M)}
\properrow{Aquae Salviae / Tre Fontane; three springs}{Pseudo-Marcellus
(5th c.)}{monastery and church at the site}{``most legendary''
(Benedict XVI); etiological (L)}
\properrow{The physical description}{\work{Acts of Paul} (c.~160)}{whole
iconographic type; paintings attested by Eusebius, HE 7.18}{literary;
possibly reminiscent; unverifiable (L)}
\properrow{Thecla}{\work{Acts of Paul} (c.~160)}{major early cult
(Seleucia)}{romance; evidence of reception, not of Paul (L)}
\properrow{Milk at the beheading; apparition to Nero}{\work{Acts of
Paul} martyrdom}{devotional retellings}{symbolic legend (L)}
\properrow{Seneca correspondence}{the forged letters; Jerome, \work{De
viris} 12 (392) attests their circulation}{medieval popularity}{forgery
by common consent (L)}
\properrow{Sarcophagus remains are Paul's}{site tradition from Gaius;
probe results announced 28 June 2009}{official shrine claim}{consistent,
not probative; ``would seem to confirm'' (M/K)}
\end{fourstable}
